% Fedimint Confidential eCash -- scheme description.
% Builds (see Makefile for the exact commands):
%   colour A4:            latexmk -pdf Fedimint_Blinded_eCash.tex
%   black & white A4:     \def\BW{1}     before \input (G2 elements get a tilde instead of a colour)
%   black & white A5:     \def\AFIVE{1}  in addition (wide left margin for clip binding)
%   A5 two-up on A4:      impose_a5_2up.tex places the A5 pages in cut-and-stack order
\newif\ifafive
\ifdefined\AFIVE \afivetrue \fi   % \def\AFIVE{1} before \input selects the A5 layout

\ifafive
\documentclass[10pt,a5paper]{article}
% wide left margin for clip binding, 2-up on A4 and cut in half
\usepackage[a5paper,left=2.4cm,right=1.2cm,top=1.6cm,bottom=1.8cm]{geometry}
\else
\documentclass[11pt,a4paper]{article}
\usepackage[margin=2.6cm]{geometry}
\fi

% layout-dependent table widths and sizes
\ifafive
  \newcommand{\wAttr}{4.6cm}\newcommand{\wSym}{5.2cm}\newcommand{\wSec}{6.8cm}
  \newcommand{\tabsize}{\small}\newcommand{\legendsize}{\footnotesize}
  \newcommand{\annot}[2]{\text{\small #2}}   % short equation annotations
\else
  \newcommand{\wAttr}{7.5cm}\newcommand{\wSym}{7.4cm}\newcommand{\wSec}{10.5cm}
  \newcommand{\tabsize}{}\newcommand{\legendsize}{\small}
  \newcommand{\annot}[2]{\text{#1}}
\fi
\usepackage{amsmath,amssymb}
\usepackage[table]{xcolor}
\usepackage{booktabs}
\usepackage{array}
\usepackage{longtable}
\usepackage{enumitem}
\usepackage{microtype}
\usepackage[most]{tcolorbox}
\newif\ifbw
\ifdefined\BW \bwtrue \fi   % \def\BW{1} before \input selects the B/W version

\ifbw
\usepackage[colorlinks=true,allcolors=black]{hyperref}
\else
\usepackage[colorlinks=true,linkcolor=black,citecolor=black,urlcolor=blue!60!black]{hyperref}
\fi

% ---------------------------------------------------------------------------
% Typing: every symbol is marked by the set it lives in.
%   colour version: scalars brown, G1 blue, G2 magenta, GT purple
%   B/W version:    scalars lowercase, G1 uppercase, G2 uppercase with a tilde
% ---------------------------------------------------------------------------
\definecolor{ScalarCol}{RGB}{170,85,0}    % K  = Z_p         (scalars)
\definecolor{GoneCol}{RGB}{0,90,175}      % G1               (blue)
\definecolor{GtwoCol}{RGB}{175,0,95}      % G2               (magenta)
\definecolor{GtCol}{RGB}{95,0,175}        % GT               (purple)

\ifbw
\newcommand{\K}[1]{#1}                      % scalar: lowercase
\newcommand{\Go}[1]{#1}                     % element of G1: uppercase
\newcommand{\Gt}[1]{\tilde{#1}}             % element of G2: uppercase with tilde
\newcommand{\GT}[1]{#1}                     % element of GT: pairing values
\colorlet{BoxFrame}{black!75}\colorlet{BoxBack}{black!4}
\colorlet{WarnFrame}{black!75}\colorlet{WarnBack}{black!8}
\else
\newcommand{\K}[1]{{\color{ScalarCol}#1}}   % scalar
\newcommand{\Go}[1]{{\color{GoneCol}#1}}    % element of G1
\newcommand{\Gt}[1]{{\color{GtwoCol}#1}}    % element of G2
\newcommand{\GT}[1]{{\color{GtCol}#1}}      % element of GT
\colorlet{BoxFrame}{blue!45!black}\colorlet{BoxBack}{blue!3}
\colorlet{WarnFrame}{orange!70!black}\colorlet{WarnBack}{orange!5}
\fi

\newcommand{\F}{\mathbb{K}}
\newcommand{\GG}[1]{\mathbb{G}_{#1}}
\newcommand{\Hs}{\mathsf{H}_{s}}
\newcommand{\Hg}{\mathsf{H}_{\GG{1}}}
\newcommand{\auth}{\mathsf{auth}}
\newcommand{\dtag}{\mathsf{tag}}
\newcommand{\vv}[1]{\vec{#1}}
\newcommand{\sample}{\stackrel{\$}{\leftarrow}}
\newcommand{\eqq}{\stackrel{?}{=}}

\newtcolorbox{whybox}[1][]{enhanced,breakable,colback=BoxBack,colframe=BoxFrame,
  boxrule=0.5pt,left=6pt,right=6pt,top=4pt,bottom=4pt,arc=2pt,
  fonttitle=\bfseries,title={#1}}
\newtcolorbox{warnbox}[1][]{enhanced,breakable,colback=WarnBack,colframe=WarnFrame,
  boxrule=0.5pt,left=6pt,right=6pt,top=4pt,bottom=4pt,arc=2pt,
  fonttitle=\bfseries,title={#1}}

\title{Fedimint Confidential eCash}
\author{Gazelle}
\date{July 2026 \\ \small revised September 2026}

\begin{document}
\maketitle

\section{Introduction}

Electronic cash (eCash) aims to provide a digital analogue of physical cash,
enabling users to perform electronic payments while preserving strong privacy
guarantees. Traditional eCash systems focus on anonymity by preventing the
issuer from linking the withdrawal of a coin to its subsequent spending through
the use of blind signatures. More recent constructions extend this notion of
privacy by providing \emph{confidential} eCash, where not only the identity of
the user remains protected, but also the value of the eCash note. In such
systems, the denomination is cryptographically hidden from all parties except
those authorized to learn it, while still allowing the correctness of
transactions to be publicly verified.

Designing a confidential blind eCash system requires satisfying several
fundamental security and practical properties. First, the system should
guarantee \textbf{unlinkability}, ensuring that an issued note cannot be linked
to the transaction in which it is later spent, thereby preserving the user's
anonymity. Second, it must satisfy \textbf{unforgeability}, meaning that every
valid eCash note must originate from the mint and that no adversary can create
additional spendable notes or increase their value. Third, the protocol should
support \textbf{offline transactions}, allowing payments to be completed
without requiring communication with the mint during the spending phase.
Finally, the construction should be \textbf{practical}, with cryptographic
proofs and transaction sizes that remain compact and computational costs that
are sufficiently low for deployment in real-world payment systems.

This document introduces the concept of confidential blind eCash generated
within a federation and discusses the cryptographic techniques that enable
both transaction privacy and value confidentiality while maintaining strong
security guarantees and practical efficiency. Section~\ref{sec:notation}
fixes the notation, Section~\ref{sec:scheme} explains the scheme for a single
note, and only afterwards Section~\ref{sec:batching} introduces the batching
of issuance requests as an optimisation on top of it.

\section{Threshold Anonymous Credentials}

This project builds upon a modified version of the Coconut credential
scheme~\cite{coconut}, as adapted and deployed by Nym
Technologies~\cite{nym,nymcode}. Coconut is a threshold anonymous credential
system that enables users to obtain credentials from multiple authorities while
preserving privacy through selective disclosure and unlinkability. The
underlying signature is the Pointcheval--Sanders (PS) signature~\cite{ps16},
whose signing operation is linear in the secret key and can therefore be
distributed with Shamir secret sharing, and whose signatures can be
re-randomised so that a credential can be shown without being recognised.

Nym's variant differs from the original Coconut construction in the issuance
protocol: instead of ElGamal-encrypting the hidden attributes for the
authorities, the user sends Pedersen commitments to them in a base derived from
a hash. This removes the user's ElGamal key pair, shrinks the request, and
makes unblinding a matter of subtracting public-key material. Nym also
publishes the verification key in both source groups of the pairing so that
unblinding can be carried out in $\GG{1}$. The following sections describe the
credential structure adopted in this project, which follows Nym's variant, and
highlight where it deviates from it.

At a high level, a credential consists of a set of cryptographic commitments to
user attributes together with a threshold signature issued jointly by a
committee of authorities. The commitments allow users to prove statements about
their attributes in zero knowledge without revealing the attributes themselves,
while the threshold signature guarantees that the credential has been
legitimately issued even if only a quorum of authorities participates.

\section{Notation and Setting}
\label{sec:notation}

\subsection{Scalars, groups and the pairing}

Throughout, $\F := \mathbb{Z}/p\mathbb{Z}$ is the field of integers modulo the
prime $p$, the order of the BLS12-381 groups. Let $\GG{1}, \GG{2}, \GG{T}$ be the
BLS12-381 pairing groups, all of order $p$ and \emph{written additively}, and let
\[
  e : \GG{1} \times \GG{2} \to \GG{T}
\]
be the pairing. It is bilinear: for all $\K{a},\K{b} \in \F$,
$\Go{U} \in \GG{1}$ and $\Gt{W} \in \GG{2}$,
\[
  e(\K{a}\,\Go{U},\; \K{b}\,\Gt{W}) \;=\; \K{ab}\; \GT{e(\Go{U}, \Gt{W})} .
\]
Since each group has prime order $p$, scalar multiplication
$\F \times \GG{i} \to \GG{i}$, $(\K{\lambda}, \Go{U}) \mapsto \K{\lambda}\,\Go{U}$
makes every $\GG{i}$ a one-dimensional $\F$-vector space. All maps in this
document that send attributes and randomness to group elements are therefore
$\F$-\emph{linear maps}, and the whole scheme can be read in the language of
linear algebra. There is no efficient map from $\GG{1}$ to $\GG{2}$ or back
(type-3 pairing), which is what makes it safe to publish key material in both
groups.

\begin{tcolorbox}[enhanced,colback=gray!4,colframe=gray!60!black,boxrule=0.5pt,
  arc=2pt,title={\bfseries \ifbw Notation legend: every symbol is typeset by the set it lives in\else Colour legend: every symbol is coloured by the set it lives in\fi}]
\legendsize
\ifafive\begin{tabular}{@{}lp{2.5cm}p{2.3cm}@{}}\else\begin{tabular}{@{}lll@{}}\fi
  $\K{m, r, x, \rho, s, \gamma, \lambda}$ & scalars in $\F$ (32\,B) & lowercase \\
  $\Go{G_1, G_p, H_p, H_j, B, C_j, P, Y_j, \sigma, R}$ & elements of $\GG{1}$ (48\,B) & uppercase or Greek \\
  $\Gt{G_2}, \Gt{X_j}, \Gt{K'}, \Gt{V}$ & elements of $\GG{2}$ (96\,B) & uppercase\ifbw{} with a tilde\fi \\
  $\GT{e(\cdot,\cdot)}$ & elements of $\GG{T}$ & only inside pairing checks \\
\end{tabular}

\medskip
Groups are written additively. Group elements are uppercase letters (or Greek,
as for the signature $\Go{\sigma}$); scalars are lowercase\ifbw; elements of $\GG{2}$
additionally carry a tilde, as in Pointcheval--Sanders~\cite{ps16}\fi. Sizes are for compressed encodings. A scalar times a $\GG{1}$ element is a
$\GG{1}$ element; a scalar times a $\GG{2}$ element is a $\GG{2}$ element;
elements of different groups are never added to each other. Vectors are written
with an arrow ($\vv{w}$, $\vv{C}$) and are typed component-wise.
\end{tcolorbox}

\subsection{Generators and hash functions}

$\Go{G_1} \in \GG{1}$ and $\Gt{G_2} \in \GG{2}$ are fixed generators. All further
generators are derived by hashing domain-separation strings to the curve
(nothing-up-my-sleeve), so that nobody knows a discrete logarithm relation
between any two of them:
\begin{itemize}[nosep]
  \item $\Go{G_p}, \Go{H_p} \in \GG{1}$ -- Pedersen basis for \emph{amounts};
  \item $\Go{H_1}, \Go{H_2}, \Go{H_3} \in \GG{1}$ -- basis for the attribute commitment $\Go{C_m}$.
\end{itemize}
The generator $\Go{G_1}$ doubles as the blinding generator of all commitments
that the mint signs. This is not a coincidence: unblinding in
Section~\ref{sec:unblind} only works because the blinding generator of the
commitments is the same generator in which the mint's public key is published
in $\GG{1}$.

Two hash functions, modelled as random oracles, are used:
\[
  \Hg : \{0,1\}^* \to \GG{1} \qquad\text{and}\qquad \Hs : \{0,1\}^* \to \F .
\]
$\Hg$ derives the per-note signature base $\Go{B}$; $\Hs$ derives
Fiat--Shamir challenges and encodes the spend condition as a scalar. $\|$
denotes concatenation of canonical encodings.

\subsection{Attributes}

Each eCash note carries three attributes in $\F$. They are the messages that
the federation signs, and they are never revealed to the mint at issuance.

\begin{center}
\tabsize
\begin{tabular}{@{}llp{\wAttr}@{}}
\toprule
Attribute & Meaning & Origin \\
\midrule
$\K{m_1}$ & the \emph{amount} & a 64-bit integer lifted to $\F$ \\
$\K{m_2}$ & the \emph{serial secret} & uniformly random; basis of the nullifier \\
$\K{m_3} = \Hs(\auth)$ & the \emph{spend condition} & hash of authentication data fixed at issuance \\
\bottomrule
\end{tabular}
\end{center}

\subsection{Keys}

For a single issuer, the secret key is a vector of four scalars,
\[
  \mathsf{sk} = (\K{x_0}, \K{x_1}, \K{x_2}, \K{x_3}) \in \F^4,
\]
one \emph{constant term} $\K{x_0}$ plus one scalar $\K{x_j}$ per attribute
$\K{m_j}$. The public key publishes each component in \emph{both} source groups:
\[
  \mathsf{pk} \;=\; \bigl(\;
    \underbrace{\{\Go{Y_j} = \K{x_j}\,\Go{G_1}\}_{j=0}^{3}}_{\text{in } \GG{1}\text{: used for unblinding}},\quad
    \underbrace{\{\Gt{X_j} = \K{x_j}\,\Gt{G_2}\}_{j=0}^{3}}_{\text{in } \GG{2}\text{: used for verification}}
  \;\bigr).
\]
In Coconut's notation $\Gt{X_0}$ is $\alpha$, $\Gt{X_j}$ is $\beta_j$ and
$\Go{Y_j}$ is Nym's $\beta_{g_1,j}$~\cite{nymcode}.
Section~\ref{sec:threshold} replaces the single issuer by a $t$-of-$n$
federation whose \emph{aggregate} key has exactly this shape.

\subsection{Summary of symbols}

Table~\ref{tab:symbols} lists every symbol used in the remainder of the
document together with the set it lives in.

\newcommand{\symbolrows}{%
$p$, $\F$ & -- & group order and scalar field $\mathbb{Z}/p\mathbb{Z}$ \\
$\Go{G_1}$, $\Gt{G_2}$ & $\GG{1}$, $\GG{2}$ & fixed generators (Coconut's $g_1$, $g_2$); $\Go{G_1}$ is also the blinding generator \\
$\Go{G_p}, \Go{H_p}$ & $\GG{1}$ & Pedersen basis for amounts \\
$\Go{H_1}, \Go{H_2}, \Go{H_3}$ & $\GG{1}$ & basis for the attribute commitment \\
$\K{m_1}, \K{m_2}, \K{m_3}$ & $\F$ & amount, serial secret, spend condition \\
$\auth$ & bytes & authentication data with $\K{m_3} = \Hs(\auth)$ \\
$\K{r_p}, \K{r_m}, \K{r_1}, \K{r_2}, \K{r_3}$ & $\F$ & commitment randomness chosen by the user at issuance \\
$\Go{P}$ & $\GG{1}$ & Pedersen commitment to the amount \\
$\Go{C_m}$ & $\GG{1}$ & commitment to all attributes; input to $\Hg$ \\
$\Go{B}$ & $\GG{1}$ & signature base, $\Go{B} = \Hg(\Go{C_m})$ (Coconut's $h$) \\
$\Go{C_1}, \Go{C_2}, \Go{C_3}$ & $\GG{1}$ & blinded attributes (commitments in base $\Go{B}$) \\
$\vv{C}$ & $\GG{1}^5$ & the statement of the issuance proof, $(\Go{P}, \Go{C_m}, \Go{C_1}, \Go{C_2}, \Go{C_3})$ \\
$\vv{w}$ & $\F^8$ & the witness of the issuance proof, $(\K{m_1},\K{m_2},\K{m_3},\K{r_p},\K{r_m},\K{r_1},\K{r_2},\K{r_3})$ \\
$\varphi_{\Go{B}}$ & $\F^8 \to \GG{1}^5$ & linear map with $\varphi_{\Go{B}}(\vv{w}) = \vv{C}$ \\
$\vv{\rho}$, $\vv{R}$, $\K{\gamma}$, $\vv{s}$ & $\F^8$, $\GG{1}^5$, $\F$, $\F^8$ & $\Sigma$-protocol nonce, commitment, challenge, response \\
$\K{x_0}, \dots, \K{x_3}$ & $\F$ & secret key \\
$\Go{Y_j} = \K{x_j}\Go{G_1}$, $\Gt{X_j} = \K{x_j}\Gt{G_2}$ & $\GG{1}$, $\GG{2}$ & public key in both groups \\
$\Go{\widetilde{\sigma}}$, $\Go{\sigma}$ & $\GG{1}$ & blinded and unblinded signature; $(\Go{B}, \Go{\sigma})$ is the credential \\
$\K{\lambda_i}$, $S$ & $\F$, set & Lagrange coefficients and the set of $t$ contributing mints \\
$\K{r}, \K{r'}, \K{r'_p}$ & $\F$ & re-randomisation and fresh commitment randomness at spend time \\
$(\Go{B'}, \Go{\sigma'}, \Gt{V})$ & $\GG{1}, \GG{1}, \GG{2}$ & re-randomised signature shown at spend time \\
$\Go{P'}$ & $\GG{1}$ & fresh Pedersen commitment to the amount at spend time \\
$\Gt{K'}$ & $\GG{2}$ & nullifier, $\Gt{K'} = \K{m_1}\Gt{X_1} + \K{m_2}\Gt{X_2}$ \\
$\psi$ & $\F^3 \to \GG{1}\times\GG{2}$ & linear map of the spend proof \\
}
\ifafive
\footnotesize
\begin{longtable}{@{}llp{\wSym}@{}}
\caption{Symbols, their types and their meaning.}\label{tab:symbols}\\
\toprule
Symbol & Type & Meaning \\
\midrule
\endfirsthead
\toprule
Symbol & Type & Meaning \\
\midrule
\endhead
\bottomrule
\endlastfoot
\symbolrows
\end{longtable}
\else
\begin{table}[htbp]
\centering
\caption{Symbols, their types and their meaning.}
\label{tab:symbols}
\small
\begin{tabular}{@{}llp{\wSym}@{}}
\toprule
Symbol & Type & Meaning \\
\midrule
\symbolrows
\bottomrule
\end{tabular}
\end{table}
\fi


\section{The Scheme for a Single Note}
\label{sec:scheme}

This section presents the lifecycle of \emph{one} eCash note, first with a
single issuer and then with a federation of mints. Batching several notes into
one request is deferred to Section~\ref{sec:batching}; nothing in this section
depends on it.

\subsection{The underlying signature}
\label{sec:ps}

A Pointcheval--Sanders signature~\cite{ps16} on the attributes
$(\K{m_1},\K{m_2},\K{m_3})$ with base point $\Go{B} \in \GG{1}$ (Coconut's $h$) is the pair
$(\Go{B}, \Go{\sigma})$ with
\begin{equation}
  \Go{\sigma} \;=\; \Bigl(\K{x_0} + \sum_{j=1}^{3} \K{x_j}\,\K{m_j}\Bigr)\,\Go{B},
  \label{eq:ps}
\end{equation}
verified with the public key in $\GG{2}$ by the single pairing equation
\begin{equation}
  e\Bigl(\Go{B},\; \Gt{X_0} + \sum_{j=1}^{3} \K{m_j}\,\Gt{X_j}\Bigr) \;\eqq\; e(\Go{\sigma},\, \Gt{G_2}).
  \label{eq:psverify}
\end{equation}
Correctness is bilinearity: both sides equal
$\GT{e(\Go{B},\Gt{G_2})}$ raised to $\K{x_0} + \sum_j \K{x_j}\K{m_j}$.
Three properties of this signature carry the whole scheme.

\begin{whybox}[Why PS signatures]
\begin{enumerate}[nosep,leftmargin=*]
\item \textbf{Signing is linear in the key.} Equation~\eqref{eq:ps} is a
  linear function of $(\K{x_0},\dots,\K{x_3})$. Anything linear in the key can
  be Shamir-shared and interpolated in the exponent, which gives threshold
  issuance (Section~\ref{sec:threshold}).
\item \textbf{Signing is linear in the messages.} A commitment $\K{m_j}\Go{B} + \K{r_j}\Go{G_1}$
  can be fed into~\eqref{eq:ps} in place of $\K{m_j}\Go{B}$; the signer never
  sees $\K{m_j}$, and the extra term is removable. This gives blind issuance
  (Section~\ref{sec:blindsign}).
\item \textbf{Re-randomisable.} For any $\K{r} \neq 0$, $(\K{r}\Go{B}, \K{r}\Go{\sigma})$
  is an equally valid signature on the same messages, and is unlinkable to
  $(\Go{B},\Go{\sigma})$ under DDH in $\GG{1}$. This gives unlinkable spending
  (Section~\ref{sec:spend}).
\end{enumerate}
\end{whybox}

The base $\Go{B}$ must be different for every signature. If two signatures share
a base, their affine combinations are again valid signatures, on the
corresponding affine combination of the messages (Coconut, footnote~3
in~\cite{coconut}; see also Section~\ref{sec:sharedh}). In plain PS the signer
picks $\Go{B}$ at random, which a set of independent mints cannot do
consistently; Coconut instead derives $\Go{B}$ by hashing a commitment supplied
by the user, so that every mint computes the same $\Go{B}$ on its own.

\subsection{Requesting issuance}
\label{sec:request}

The user holds the attributes $(\K{m_1},\K{m_2},\K{m_3})$, samples the
randomness $\K{r_p}, \K{r_m}, \K{r_1}, \K{r_2}, \K{r_3}$ uniformly from $\F$,
and forms five elements of $\GG{1}$:
\begin{align}
  \Go{P} &= \K{m_1}\,\Go{G_p} + \K{r_p}\,\Go{H_p}
     && \annot{(Pedersen commitment to the amount)}{(amount commitment)} \label{eq:pc} \\
  \Go{C_m} &= \K{m_1}\,\Go{H_1} + \K{m_2}\,\Go{H_2} + \K{m_3}\,\Go{H_3} + \K{r_m}\,\Go{G_1}
     && \annot{(commitment to all attributes)}{(attribute commitment)} \label{eq:cm} \\
  \Go{B} &= \Hg(\Go{C_m})
     && \annot{(derived signature base)}{(signature base)} \label{eq:h} \\
  \Go{C_j} &= \K{m_j}\,\Go{B} + \K{r_j}\,\Go{G_1}
     && \annot{(blinded attributes, $j \in \{1,2,3\}$)}{(blinded attr., $j \in \{1,2,3\}$)} \label{eq:cj}
\end{align}
The issuance request consists of the statement
$\vv{C} = (\Go{P}, \Go{C_m}, \Go{C_1}, \Go{C_2}, \Go{C_3}) \in \GG{1}^5$, the
proof of Section~\ref{sec:issuanceproof}, and the range proof of
Section~\ref{sec:range}. The mint recomputes $\Go{B}$ from $\Go{C_m}$ itself;
it is not transmitted.

\begin{whybox}[What each element is for]
\begin{description}[nosep,leftmargin=1.2em]
\item[$\Go{P}$ -- the amount commitment.] Pedersen commitment in its own
  basis $(\Go{G_p},\Go{H_p})$: perfectly hiding because of $\K{r_p}$, binding
  under the discrete-logarithm assumption between $\Go{G_p}$ and $\Go{H_p}$. It is
  \emph{additively homomorphic}, $\Go{P}(\K{m},\K{r}) + \Go{P}(\K{m'},\K{r'}) = \Go{P}(\K{m+m'},\K{r+r'})$,
  which lets the transaction layer check that inputs and outputs balance
  without seeing the amounts. It is also the object the range proof
  (Section~\ref{sec:range}) is attached to. It is deliberately in a basis
  unrelated to the signature so that balance checking is independent of the
  credential.
\item[$\Go{C_m}$ -- the commitment behind $\Go{B}$.] A Pedersen vector
  commitment to all three attributes with randomness $\K{r_m}$. Its only
  purpose is to be hashed into $\Go{B}$. Because it is \emph{binding},
  $\Go{B}$ is a function of the attributes the user is about to get signed;
  because it is \emph{hiding}, the mint learns nothing from it; because
  $\K{r_m}$ is fresh, every request gets a fresh base even for identical
  attributes.
\item[$\Go{B}$ -- the signature base.] Neither party chooses it: the user
  cannot pick a base with a known relation to another note's base (which would
  enable the affine-combination forgery of Section~\ref{sec:ps}), and the mint
  does not have to pick anything, so all mints of a federation agree on
  $\Go{B}$ without communicating. This is the Coconut trick that turns PS into
  a threshold-friendly scheme.
\item[$\Go{C_j}$ -- the blinded attributes.] Pedersen commitments to each
  attribute \emph{in the base} $\Go{B}$ with blinding on $\Go{G_1}$. The base
  $\Go{B}$ is what the signature~\eqref{eq:ps} multiplies the attributes with,
  so the mint can apply its key directly to $\Go{C_j}$
  (Section~\ref{sec:blindsign}). The blinding generator $\Go{G_1}$ is what the
  public key $\Go{Y_j}$ is published in, so the blinding term the mint
  inadvertently signs, $\K{r_j}\K{x_j}\Go{G_1} = \K{r_j}\Go{Y_j}$, is computable by the user
  from public data. In the original Coconut these were ElGamal ciphertexts of
  $\K{m_j}\Go{B}$; Nym's commitments serve the same purpose with less data.
\end{description}
\end{whybox}

\subsection{The issuance proof}
\label{sec:issuanceproof}

The mint must be convinced that $\vv{C}$ is well formed before signing:
otherwise a user could, for instance, put one amount into $\Go{P}$ (and its
range proof) and a different amount into $\Go{C_1}$ (and the signature). Well-formedness
of $\vv{C}$ is exactly the statement that $\vv{C}$ lies in the image of the
linear map $\varphi_{\Go{B}} : \F^8 \to \GG{1}^5$,
\begin{equation}
  \ifafive \varphi_{\Go{B}}(\vv{w}) \else
  \varphi_{\Go{B}}(\K{m_1},\K{m_2},\K{m_3},\K{r_p},\K{r_m},\K{r_1},\K{r_2},\K{r_3}) \fi
  \;=\;
  \begin{pmatrix}
    \K{m_1}\,\Go{G_p} + \K{r_p}\,\Go{H_p} \\[2pt]
    \K{m_1}\,\Go{H_1} + \K{m_2}\,\Go{H_2} + \K{m_3}\,\Go{H_3} + \K{r_m}\,\Go{G_1} \\[2pt]
    \K{m_1}\,\Go{B} + \K{r_1}\,\Go{G_1} \\[2pt]
    \K{m_2}\,\Go{B} + \K{r_2}\,\Go{G_1} \\[2pt]
    \K{m_3}\,\Go{B} + \K{r_3}\,\Go{G_1}
  \end{pmatrix}
  \;=\;
  \begin{pmatrix} \Go{P} \\[2pt] \Go{C_m} \\[2pt] \Go{C_1} \\[2pt] \Go{C_2} \\[2pt] \Go{C_3} \end{pmatrix}
  = \vv{C}.
  \label{eq:phi}
\end{equation}
The witness is
\[
  \vv{w} = (\K{m_1},\K{m_2},\K{m_3},\K{r_p},\K{r_m},\K{r_1},\K{r_2},\K{r_3}) \in \F^8 .
\]
Note that $\varphi_{\Go{B}}$ depends on $\Go{B}$, which the verifier recomputes
from the second component of $\vv{C}$; the map is fixed once $\Go{C_m}$ is.

The user proves knowledge of a preimage with the standard $\Sigma$-protocol
for linear maps (Schnorr's protocol generalised to a vector-valued linear map),
made non-interactive with the Fiat--Shamir transform under the domain tag
$\dtag_{\mathrm{iss}} = {}$\texttt{FEDIMINT\_\allowbreak ECASH\_\allowbreak CHALLENGE\_\allowbreak ISSUANCE}:
\begin{equation}
\begin{aligned}
  \vv{\rho} &\sample \F^8, &
  \vv{R} &= \varphi_{\Go{B}}(\vv{\rho}) \in \GG{1}^5, \\
  \K{\gamma} &= \Hs\bigl(\dtag_{\mathrm{iss}} \,\|\, \vv{C} \,\|\, \vv{R}\bigr) \in \F, &
  \vv{s} &= \vv{\rho} + \K{\gamma}\,\vv{w} \in \F^8 .
\end{aligned}
  \label{eq:sigma}
\end{equation}
The proof is $(\vv{R}, \vv{s})$: five $\GG{1}$ elements and eight scalars,
$5 \cdot 48 + 8 \cdot 32 = 496$ bytes. The verifier recomputes $\Go{B}$ and
$\K{\gamma}$ and checks
\begin{equation}
  \varphi_{\Go{B}}(\vv{s}) \;\eqq\; \K{\gamma}\,\vv{C} + \vv{R}.
  \label{eq:sigmaverify}
\end{equation}

\begin{whybox}[Why this proves what we need]
\begin{description}[nosep,leftmargin=1.2em]
\item[Completeness.] By linearity,
  $\varphi_{\Go{B}}(\vv{s}) = \varphi_{\Go{B}}(\vv{\rho}) + \K{\gamma}\,\varphi_{\Go{B}}(\vv{w}) = \vv{R} + \K{\gamma}\vv{C}$.
\item[Knowledge (special soundness).] Two accepting transcripts with the same
  $\vv{R}$ and different challenges $\K{\gamma} \neq \K{\gamma'}$ give
  $\varphi_{\Go{B}}(\vv{s}-\vv{s}') = (\K{\gamma}-\K{\gamma'})\,\vv{C}$, so
  $(\K{\gamma}-\K{\gamma'})^{-1}(\vv{s}-\vv{s}')$ is a witness. A prover who can
  answer two challenges knows a preimage. In the random-oracle model
  Fiat--Shamir preserves this.
\item[Zero knowledge.] A simulator picks $\vv{s} \sample \F^8$ and
  $\K{\gamma} \sample \F$ first and sets $\vv{R} = \varphi_{\Go{B}}(\vv{s}) - \K{\gamma}\vv{C}$;
  the transcript is distributed exactly like a real one. The mint learns
  nothing about the attributes beyond the commitments, which are hiding.
\item[What the single witness buys us.] One witness vector serves all five
  equations, so the \emph{same} $\K{m_1}$ sits in $\Go{P}$, $\Go{C_m}$ and $\Go{C_1}$,
  and the same $\K{m_2},\K{m_3}$ sit in $\Go{C_m}$ and $\Go{C_2},\Go{C_3}$. Hence
  (i) the amount that is range-proven is the amount that gets signed
  (unforgeability of value), (ii) $\Go{B}$ really is derived from a commitment
  to the attributes being signed (the freshness argument for $\Go{B}$ holds),
  and (iii) the $\Go{C_j}$ are of the exact form~\eqref{eq:cj}, so the
  blinding term in the signature is exactly $\sum_j \K{r_j}\Go{Y_j}$ and
  nothing else gets signed. Hashing $\vv{C}$ into the challenge binds the proof
  to the statement; the domain tag prevents replaying a proof for $\psi$ of
  Section~\ref{sec:spend} against $\varphi_{\Go{B}}$.
\end{description}
\end{whybox}

\subsection{Range proof on the amount}
\label{sec:range}

The commitment $\Go{P}$ hides $\K{m_1}$ as an element of $\F$, and $\F$ wraps
around: a ``negative'' amount $-\K{v}$ is the field element $p - \K{v}$. Without
further constraints an adversary could commit to an out-of-range value and
satisfy a balance equation $\sum \Go{P}_{\mathrm{in}} = \sum \Go{P}_{\mathrm{out}}$
while creating value from nothing. The user therefore additionally proves that
the committed amount lies in an accepted interval, e.g.\ $0 \le \K{m_1} < 2^{64}$.
This work employs Bulletproofs~\cite{bulletproofs}, which give short
non-interactive range proofs over Pedersen commitments without a trusted setup,
with proof size logarithmic in the bit length and aggregatable across several
commitments. Because the issuance proof ties the $\K{m_1}$ in $\Go{P}$ to the
$\K{m_1}$ in the signature, the range proof on $\Go{P}$ carries over to the
signed amount.

\subsection{Blind signing}
\label{sec:blindsign}

A mint that has verified the proofs recomputes $\Go{B} = \Hg(\Go{C_m})$ and
applies its key to the request as if the $\Go{C_j}$ were the attributes:
\begin{equation}
  \Go{\widetilde{\sigma}}
  \;=\; \K{x_0}\,\Go{B} + \sum_{j=1}^{3} \K{x_j}\,\Go{C_j}
  \;=\; \underbrace{\Bigl(\K{x_0} + \sum_{j=1}^{3} \K{x_j}\,\K{m_j}\Bigr)\Go{B}}_{\text{PS signature } \Go{\sigma}\text{, eq.~\eqref{eq:ps}}}
      \;+\; \underbrace{\sum_{j=1}^{3} \K{r_j}\,\bigl(\K{x_j}\,\Go{G_1}\bigr)}_{\text{blinding term } = \sum_j \K{r_j}\Go{Y_j}} .
  \label{eq:blindsign}
\end{equation}
The expansion is just~\eqref{eq:cj} substituted into the first expression. The
mint never sees $\K{m_j}$; it sees $\Go{B}$, the $\Go{C_j}$ and the proof, all
of which are independent of the attributes by the hiding property and zero
knowledge.

\subsection{Unblinding and verification}
\label{sec:unblind}

The blinding term in~\eqref{eq:blindsign} is a linear combination of the
$\GG{1}$ public key components with coefficients the user chose, so the user
removes it from public information only:
\begin{equation}
  \Go{\sigma} \;=\; \Go{\widetilde{\sigma}} - \sum_{j=1}^{3} \K{r_j}\,\Go{Y_j},
  \label{eq:unblind}
\end{equation}
and then checks $(\Go{B},\Go{\sigma})$ against the $\GG{2}$ key with the PS
equation~\eqref{eq:psverify}. This is where the choice of $\Go{G_1}$ as the
blinding generator in~\eqref{eq:cj} is essential: with any other blinding
generator the user would need $\K{x_j}$ to cancel the term. The pair
$(\Go{B},\Go{\sigma})$ together with $(\K{m_1},\K{m_2},\auth)$ is the eCash note.

\subsection{Threshold issuance}
\label{sec:threshold}

Nothing above requires a single issuer, because~\eqref{eq:blindsign} is
$\F$-linear in the secret key $(\K{x_0},\dots,\K{x_3})$. The standard
threshold-BLS transformation therefore applies component-wise.

\paragraph{Key sharing.}
Each $\K{x_j}$ is Shamir-shared: pick random polynomials $f_j \in \F[X]$ of
degree at most $t-1$ with $f_j(0) = \K{x_j}$, and give mint
$i \in \{1,\dots,n\}$ the key share
\[
  \mathsf{sk}_i = \bigl(\K{f_0(i)}, \K{f_1(i)}, \K{f_2(i)}, \K{f_3(i)}\bigr) \in \F^4,
\]
with public key shares $\Go{Y_{i,j}} = \K{f_j(i)}\,\Go{G_1}$ and
$\Gt{X_{i,j}} = \K{f_j(i)}\,\Gt{G_2}$. The aggregate public key is
$\Go{Y_j} = \K{f_j(0)}\,\Go{G_1}$, $\Gt{X_j} = \K{f_j(0)}\,\Gt{G_2}$, i.e.\ exactly
the single-issuer key of Section~\ref{sec:notation}. Fewer than $t$ shares
reveal nothing about $\K{x_j}$.

\paragraph{Share signing and share verification.}
Every mint runs Section~\ref{sec:blindsign} with its own share,
\[
  \Go{\widetilde{\sigma}_i} \;=\; \K{f_0(i)}\,\Go{B} + \sum_{j=1}^{3} \K{f_j(i)}\,\Go{C_j},
\]
and the user unblinds each share with the mint's $\GG{1}$ key share,
$\Go{\sigma_i} = \Go{\widetilde{\sigma}_i} - \sum_j \K{r_j}\,\Go{Y_{i,j}}$, then
verifies it with~\eqref{eq:psverify} against $\Gt{X_{i,j}}$. A mint that
returns garbage is identified immediately and simply left out.

\paragraph{Aggregation.}
Let $S$ be a set of $t$ mint indices whose shares verified. The federation
signature is the Lagrange combination at $0$,
\begin{equation}
  \Go{\sigma} \;=\; \sum_{i \in S} \K{\lambda_i}\,\Go{\sigma_i},
  \qquad
  \K{\lambda_i} \;=\; \prod_{\substack{k \in S \\ k \neq i}} \frac{k}{k - i} \pmod p .
  \label{eq:lagrange}
\end{equation}
Since $\Go{\sigma_i} = \bigl(f_0(i) + \sum_j \K{m_j} f_j(i)\bigr)\,\Go{B}$ and
$\K{\lambda} \mapsto \K{\lambda}\,\Go{B}$ is linear, the sum evaluates the
polynomial $f_0 + \sum_j \K{m_j} f_j$ at $0$ inside the group, which is exactly
the single-issuer signature~\eqref{eq:ps}. Everything downstream is unchanged
and verifies against the aggregate key.

\subsection{Spending}
\label{sec:spend}

To spend, the user must show the mint a valid signature on a note without
letting the mint connect it to the issuance it came from, while still letting
the mint (i) check the amount against the transaction balance and (ii) detect
a second spend of the same note.

\paragraph{Re-randomisation.}
The user samples $\K{r}, \K{r'} \sample \F$ and forms
\begin{equation}
  (\Go{B'},\; \Go{\sigma'},\; \Gt{V}) \;=\; \bigl(\K{r}\,\Go{B},\;\; \K{r}\,\Go{\sigma} + \K{r}\K{r'}\,\Go{B},\;\; \K{r'}\,\Gt{G_2}\bigr).
  \label{eq:rerand}
\end{equation}
The pair $(\K{r}\Go{B}, \K{r}\Go{\sigma})$ is the PS re-randomisation of
Section~\ref{sec:ps}: fresh group elements, unlinkable to $(\Go{B},\Go{\sigma})$
under DDH in $\GG{1}$. The $\K{r'}$ terms are inherited from Coconut's
$(\kappa,\nu)$ and are discussed in Section~\ref{sec:V}.

\paragraph{Shown elements.}
Along with $(\Go{B'},\Go{\sigma'},\Gt{V})$ and the authentication data
$\auth$ the user sends
\begin{equation}
  \Go{P'} = \K{m_1}\,\Go{G_p} + \K{r'_p}\,\Go{H_p} \in \GG{1},
  \qquad
  \Gt{K'} = \K{m_1}\,\Gt{X_1} + \K{m_2}\,\Gt{X_2} \in \GG{2},
  \label{eq:spendelems}
\end{equation}
where $\K{r'_p} \sample \F$ is \emph{fresh} commitment randomness, together with a
second Fiat--Shamir $\Sigma$-proof (tag
$\dtag_{\mathrm{sp}} = {}$\texttt{FEDIMINT\_\allowbreak ECASH\_\allowbreak CHALLENGE\_\allowbreak SPEND}) of knowledge of
$(\K{m_1},\K{m_2},\K{r'_p})$ for the linear map
\begin{equation}
\begin{aligned}
  \psi &: \F^3 \to \GG{1} \times \GG{2}, \\
  \psi(\K{m_1}, \K{m_2}, \K{r'_p}) &= \bigl(\, \K{m_1}\,\Go{G_p} + \K{r'_p}\,\Go{H_p},\;\; \K{m_1}\,\Gt{X_1} + \K{m_2}\,\Gt{X_2} \,\bigr) = (\Go{P'}, \Gt{K'}).
\end{aligned}
  \label{eq:psi}
\end{equation}
The protocol is the same as in~\eqref{eq:sigma} with $\psi$ in place of
$\varphi_{\Go{B}}$ and witness $(\K{m_1},\K{m_2},\K{r'_p})$; the prototype
transmits the challenge instead of the commitment $(\Go{R_{P'}},\Gt{R_K})$, so a
spend proof is four scalars ($128$ bytes), and the verifier recomputes the
commitment as $\psi(\vv{s}) - \K{\gamma}\,(\Go{P'},\Gt{K'})$.

\paragraph{Verification.}
The mint recomputes $\K{m_3} = \Hs(\auth)$, requires $\K{m_3} \neq 0$ and
$\Go{B'} \neq 0$, verifies the $\psi$-proof, checks that $\Gt{K'}$ has not been
seen before, and finally checks the signature:
\begin{equation}
  e\bigl(\Go{B'},\; \Gt{X_0} + \Gt{K'} + \K{m_3}\,\Gt{X_3} + \Gt{V}\bigr) \;\eqq\; e(\Go{\sigma'},\, \Gt{G_2}).
  \label{eq:spendverify}
\end{equation}
Expanding both sides shows why it holds. Write
$\K{z} := \K{x_0} + \K{x_1}\K{m_1} + \K{x_2}\K{m_2} + \K{x_3}\K{m_3}$, so that
$\Go{\sigma} = \K{z}\,\Go{B}$ and
$\Gt{X_0} + \Gt{K'} + \K{m_3}\Gt{X_3} = \K{z}\,\Gt{G_2}$. Then
\begin{align*}
  \text{LHS} &= e\bigl(\K{r}\Go{B},\; (\K{z} + \K{r'})\,\Gt{G_2}\bigr)
             = \K{r}(\K{z}+\K{r'})\,\GT{e(\Go{B},\Gt{G_2})}, \\
  \text{RHS} &= e\bigl(\K{r}\K{z}\,\Go{B} + \K{r}\K{r'}\,\Go{B},\; \Gt{G_2}\bigr)
             = \K{r}(\K{z}+\K{r'})\,\GT{e(\Go{B},\Gt{G_2})}.
\end{align*}
The equation is the PS verification~\eqref{eq:psverify} with the hidden part
$\K{m_1}\Gt{X_1} + \K{m_2}\Gt{X_2}$ of the message supplied as the single
element $\Gt{K'}$, the revealed part $\K{m_3}\Gt{X_3}$ recomputed by the mint,
and the $\K{r'}$ terms cancelling on both sides.

\begin{whybox}[What each shown element achieves]
\begin{description}[nosep,leftmargin=1.2em]
\item[$(\Go{B'},\Go{\sigma'})$ -- unforgeability, unlinkability.] A valid
  re-randomised PS signature proves the federation signed \emph{some}
  attributes; fresh randomness makes the pair independent of the issuance
  transcript.
\item[$\Gt{K'}$ -- the nullifier and the hidden message.] It supplies the
  hidden attributes to the verification equation without revealing them, and
  it is \emph{deterministic} in $(\K{m_1},\K{m_2})$: spending the same note
  twice produces the same $\Gt{K'}$, so the mints keep a set of seen $\Gt{K'}$
  and reject repeats. The random $\K{m_2}$ makes $\Gt{K'}$ unique per note and,
  since $\K{m_2}$ was only ever seen blinded, unlinkable to issuance. Its
  soundness rests on the $\psi$-proof: a spender must \emph{know}
  $(\K{m_1},\K{m_2})$ behind $\Gt{K'}$, so $\Gt{K'}$ cannot be re-randomised into
  a fresh nullifier for the same note.
\item[$\Go{P'}$ -- confidential amount at spend time.] A fresh Pedersen
  commitment to the same amount, in the balance-check basis, with new
  randomness so it is unlinkable to $\Go{P}$. The $\psi$-proof ties the
  $\K{m_1}$ in $\Go{P'}$ to the $\K{m_1}$ in $\Gt{K'}$, and~\eqref{eq:spendverify}
  ties the $\K{m_1}$ in $\Gt{K'}$ to the signature. So the amount the
  transaction layer balances with is the amount the federation signed, which is
  the amount that was range-proven at issuance.
\item[$\K{m_3}$, $\auth$ -- the spend condition.] $\K{m_3}$ is baked into
  $\Go{\sigma}$, so only a spender who can present the committed $\auth$ makes
  \eqref{eq:spendverify} hold. Revealing $\auth$ does not link the spend to
  issuance because the mint only saw $\K{m_3}$ blinded. Requiring
  $\K{m_3} \neq 0$ excludes the degenerate case in which the third attribute
  contributes nothing to the signature.
\end{description}
\end{whybox}

\subsubsection{On the role of \texorpdfstring{$\Gt{V}$}{V}}
\label{sec:V}

In Coconut the spender shows $\kappa = \Gt{X_0} + \sum_j \K{m_j}\Gt{X_j} + \K{r'}\Gt{G_2}$
and $\nu = \K{r'}\Go{B'}$, and the verifier checks
$e(\Go{B'},\kappa) = e(\Go{\sigma'} + \nu, \Gt{G_2})$. The fresh $\K{r'}$ in
$\kappa$ is what hides the attributes there. In this scheme $\Gt{V} = \K{r'}\Gt{G_2}$
plays the part of $\kappa$'s blinding term and $\K{r}\K{r'}\Go{B}$ inside
$\Go{\sigma'}$ plays the part of $\nu$. However, $\Gt{K'}$ is deliberately
\emph{not} blinded, because it has to be a deterministic nullifier; and
$e(\Go{B'},\Gt{V}) = e(\K{r}\K{r'}\Go{B},\Gt{G_2})$ is computable by anyone from
the shown data, so an observer can cancel the $\K{r'}$ contribution from
$e(\Go{\sigma'},\Gt{G_2})$ and recover the same $\GG{T}$ values as if
$\K{r'} = 0$. Under this reading $\Gt{V}$ adds $96$ bytes and one scalar
multiplication but no additional hiding; the unlinkability of a spend rests on
the re-randomisation by $\K{r}$ and on the fresh $\K{r'_p}$. Whether $\Gt{V}$
should be kept is an open question (Section~\ref{sec:future}).

\subsection{Security properties}
\label{sec:security}

\begin{center}
\tabsize
\begin{tabular}{@{}lp{\wSec}@{}}
\toprule
Property & Where it comes from \\
\midrule
Unforgeability & Existential unforgeability of PS signatures (LRSW-type assumption~\cite{ps16}) with a per-note base $\Go{B} = \Hg(\Go{C_m})$; binding of $\Go{P}$, $\Go{C_m}$, $\Go{C_j}$; soundness of the $\varphi_{\Go{B}}$-proof (same $\K{m_1}$ everywhere) and of the $\psi$-proof; the range proof on $\Go{P}$. Fewer than $t$ mints cannot sign. \\
Blindness & The mint sees only hiding commitments $\vv{C}$ and zero-knowledge proofs; the blinded signature reveals nothing further. \\
Unlinkability & Re-randomisation $(\K{r}\Go{B},\K{r}\Go{\sigma})$ (DDH in $\GG{1}$) and fresh $\K{r'_p}$ in $\Go{P'}$ make the spend transcript independent of the issuance transcript; $\K{m_2}$ makes $\Gt{K'}$ independent of anything shown at issuance. \\
Double-spend detection & Equality of nullifiers $\Gt{K'}$, with the $\psi$-proof preventing a spender from producing a second nullifier for the same $(\K{m_1},\K{m_2})$. \\
Value confidentiality & $\Go{P}$ and $\Go{P'}$ are perfectly hiding; $\Gt{K'}$ hides $\K{m_1}$ under DDH in $\GG{2}$ because of the random $\K{m_2}$. \\
Offline spending & Verification~\eqref{eq:spendverify} needs only the aggregate public key; only double-spend detection needs the mints' shared state. \\
\bottomrule
\end{tabular}
\end{center}

\section{Batching Issuance Requests}
\label{sec:batching}

A naive approach to requesting several notes at once repeats
Section~\ref{sec:request} per note. The proof alone is then $496$ bytes per
note, and together with Fedimint's consensus message limit of $64$\,KB this
severely restricts how many notes a single transaction can issue. Several
techniques exist to combine multiple $\Sigma$-protocol instances into one
compact proof: folding as used by Bulletproofs for non-linear relations, or
amortisation for parallel instances of the same linear relation. The latter is
used here. There are still open discussions concerning the security of the
protocol when batching is applied, both for the proofs and for the credential
scheme itself; see Section~\ref{sec:sharedh}.

\subsection{Batched issuance proof}

Suppose the user requests $n$ notes, $i = 1,\dots,n$, with witnesses
$\vv{w}_i \in \F^8$ and statements
\[
  \vv{C}_i = (\Go{P_i}, \Go{C_{m,i}}, \Go{C_{1,i}}, \Go{C_{2,i}}, \Go{C_{3,i}}) \in \GG{1}^5 .
\]
The amortised protocol sums $n$ instances of~\eqref{eq:sigma}
into a single response, which only works if all instances use the \emph{same}
linear map. Since $\varphi_{\Go{B}}$ depends on $\Go{B}$, the batch shares one
base derived from all attribute commitments of the batch:
\begin{equation}
  \Go{B} \;=\; \Hg\bigl(\Go{C_{m,1}} \,\|\, \Go{C_{m,2}} \,\|\, \cdots \,\|\, \Go{C_{m,n}}\bigr),
  \qquad
  \varphi_{\Go{B}}(\vv{w}_i) = \vv{C}_i \quad \text{for all } i .
  \label{eq:batchh}
\end{equation}
For $n = 1$ this coincides with~\eqref{eq:h}. The prover samples one nonce
vector per instance, commits to their sum, derives one challenge \emph{per
instance}, and sums the responses:
\begin{align}
  \vv{\rho}_i &\sample \F^8, \qquad \vv{R}_i = \varphi_{\Go{B}}(\vv{\rho}_i), \qquad
  \vv{R} = \sum_{i=1}^{n} \vv{R}_i \;\in \GG{1}^5, \label{eq:batchR} \\
  \K{\gamma_i} &= \Hs\bigl(\dtag_{\mathrm{iss}} \,\|\, \vv{C}_1 \,\|\, \cdots \,\|\, \vv{C}_n \,\|\, \vv{R} \,\|\, i\bigr) \;\in \F, \label{eq:batchc} \\
  \vv{s} &= \sum_{i=1}^{n} \bigl(\vv{\rho}_i + \K{\gamma_i}\,\vv{w}_i\bigr) \;\in \F^8 . \label{eq:batchs}
\end{align}
The prover sends $(\vv{R}, \vv{s})$, again $496$ bytes regardless of $n$, along
with the statements $\vv{C}_1,\dots,\vv{C}_n$. The verifier recomputes
$\Go{B}$ from~\eqref{eq:batchh} and the $\K{\gamma_i}$ from~\eqref{eq:batchc} and
checks
\begin{equation}
  \varphi_{\Go{B}}(\vv{s}) \;\eqq\; \sum_{i=1}^{n} \K{\gamma_i}\,\vv{C}_i + \vv{R}.
  \label{eq:batchverify}
\end{equation}

\begin{whybox}[Why the batched proof works]
\begin{description}[nosep,leftmargin=1.2em]
\item[Completeness.] Linearity of $\varphi_{\Go{B}}$ applied to~\eqref{eq:batchs}:
  $\varphi_{\Go{B}}(\vv{s}) = \sum_i \varphi_{\Go{B}}(\vv{\rho}_i) + \sum_i \K{\gamma_i}\,\varphi_{\Go{B}}(\vv{w}_i) = \vv{R} + \sum_i \K{\gamma_i}\vv{C}_i$.
  This step is exactly where all instances must share $\varphi_{\Go{B}}$, i.e.\
  share $\Go{B}$.
\item[Soundness.] Equation~\eqref{eq:batchverify} is the single-instance
  check~\eqref{eq:sigmaverify} for the \emph{random linear combination}
  $\sum_i \K{\gamma_i}\vv{C}_i$ of the statements. Because every $\K{\gamma_i}$ is
  derived after all $\vv{C}_i$ and $\vv{R}$ are fixed, and the $\K{\gamma_i}$
  are distinct, a prover who does not know a preimage of some $\vv{C}_i$
  can only pass with negligible probability; this is the standard
  small-exponent/random-combination argument for batch verification of
  linear relations. Including the index $i$ in the hash gives each instance
  its own challenge, which prevents a witness of one instance from being
  moved to another.
\item[Zero knowledge.] Unchanged: $\vv{s}$ is a sum of uniformly random
  $\vv{\rho}_i$ plus witness terms, hence uniform.
\end{description}
\end{whybox}

Range proofs batch independently: Bulletproofs aggregate the $n$ commitments
$\Go{P_1},\dots,\Go{P_n}$ into one proof of size logarithmic in $n$.

\subsection{Blind signing a batch}

Each mint verifies the batched proof once and then signs every note of the
batch with the \emph{shared} base,
$\Go{\widetilde{\sigma}_i} = \K{x_0}\,\Go{B} + \sum_j \K{x_j}\,\Go{C_{j,i}}$,
exactly as in~\eqref{eq:blindsign}; unblinding, share verification,
aggregation and spending are per note and unchanged from
Section~\ref{sec:scheme}. Note that after Section~\ref{sec:spend} the notes of
one batch are shown with independent re-randomisations $\K{r}\Go{B}$, so the
shared base is not visible at spend time.

\subsection{The cost of a shared base}
\label{sec:sharedh}

Sharing $\Go{B}$ across a batch is what makes~\eqref{eq:batchverify} possible,
but it removes the guarantee of Section~\ref{sec:ps} that every signature has
its own base. All notes of a batch are PS signatures on the same $\Go{B}$, and
PS signatures on a common base are closed under affine combination: for any
scalars $\K{\lambda_i}$ with $\sum_i \K{\lambda_i} = 1$,
\[
  \sum_{i} \K{\lambda_i}\,\Go{\sigma_i}
  \;=\; \Bigl(\K{x_0} + \sum_{j=1}^{3} \K{x_j} \sum_i \K{\lambda_i}\,\K{m_{j,i}}\Bigr)\,\Go{B}
\]
is a valid signature on the attribute vector $\sum_i \K{\lambda_i}\,\vv{m}_i$.
This is the forgery Coconut's footnote~3 warns about~\cite{coconut}. For the
present scheme it means that a user holding two notes of one batch with the
same spend condition ($\K{m_{3,1}} = \K{m_{3,2}}$) and different amounts can
derive a signature on an arbitrary amount $\K{\lambda}\K{m_{1,1}} + (1-\K{\lambda})\K{m_{1,2}}$
with a fresh serial secret, and the range proofs of the original notes say
nothing about it. Whether this is exploitable end to end depends on the
transaction layer (e.g.\ whether amounts are range-proven again when spent), and
the analysis, as well as possible mitigations such as per-note bases with
per-instance responses for the $\Go{B}$-dependent components of
$\varphi_{\Go{B}}$, is listed in Section~\ref{sec:future}. Until it is
settled, the batched issuance should be regarded as a proof-size optimisation
whose security is not yet established.

\section{Benchmarks}

This section evaluates the performance of the proposed confidential federated
eCash protocol through a series of benchmarks covering its main operations.
The analysis considers the computational cost incurred by both the prover and
the verifier, the size of the generated proofs, and the storage requirements of
an eCash note throughout its lifecycle. Additionally, implementation
measurements are presented for each protocol phase to assess the protocol's
practicality and identify its computational bottlenecks.

\subsection{Proof sizes}

The proof sizes differ significantly between the issuance and spending phases
due to the different security requirements of each operation. During issuance,
the user must prove the correctness of several committed values while
preserving their confidentiality. When multiple eCash notes are requested
simultaneously, these proofs are batched as in Section~\ref{sec:batching},
which keeps the $\Sigma$-proof at a constant $496$ bytes independent of the
number of notes; only the statements $\vv{C}_i$ grow linearly.

In contrast, the spending phase only requires the user to demonstrate
ownership of a valid eCash note. This is the short $\Sigma$-proof for $\psi$,
consisting of the challenge and three responses, so the spending proof is
compact but still grows linearly with the number of notes being spent.
Table~\ref{tab:sizes} shows the proof sizes for each phase, with $\GG{1}$ and
$\GG{2}$ denoting group elements of the respective groups and $s$ a $32$-byte
scalar. For an example transaction with $100$ eCash notes as both inputs and
outputs, the complete transaction (proofs plus the statements and shown
elements of every note) takes approximately $72$\,KB, only slightly over
Fedimint's $64$\,KB consensus limit.

\begin{table}[h]
\centering
\caption{Proof size as a function of the number of eCash notes $n$. The range
proofs are assumed to be for a 32-bit range.}
\label{tab:sizes}
\tabsize
\begin{tabular}{@{}lll@{}}
\toprule
Phase & Proof size & Growth \\
\midrule
Issuance (batched $\Sigma$-proof) & $5\,\GG{1} + 8\,s$ & $O(1)$ \\
Issuance (batched range proofs) & $2\log_2(32 \times n)\,\GG{1} + 9\,\GG{1} + 2\,s$ & $O(\log_2 n)$ \\
Spending & $4 \times n \times s$ & $O(n)$ \\
\bottomrule
\end{tabular}
\end{table}

\section{Implementation Notes}
\label{sec:impl}

The prototype in this repository implements Sections~\ref{sec:scheme}
and~\ref{sec:batching} over BLS12-381, with the following deviations from the
description that are worth knowing about.

\begin{warnbox}[Prototype caveats]
\begin{itemize}[nosep,leftmargin=*]
\item \textbf{Spend-time commitment randomness.} The prototype reuses the
  issuance randomness $\K{r_p}$ when building $\Go{P'}$, so $\Go{P'} = \Go{P}$
  and a spend is trivially linkable to its issuance. The scheme requires a
  fresh $\K{r'_p}$ as in~\eqref{eq:spendelems}.
\item \textbf{Hash-to-group.} $\Hg$, $\Hs$ and the generators are implemented
  by seeding a ChaCha RNG with SHA-256 rather than a standard
  \texttt{hash\_to\_curve} ciphersuite. This is fine for
  unknown-discrete-logarithm generation but not interoperable.
\item \textbf{Checks not implemented.} $\K{m_3} \neq 0$, $\Go{B'} \neq 0$ and
  the nullifier set are not enforced by the prototype; the range proof is not
  integrated.
\item \textbf{Only the batched path is exercised.} Single-note issuance is the
  $n = 1$ case of Section~\ref{sec:batching}, and the tests run batches of two.
\end{itemize}
\end{warnbox}

\section{Future Work}
\label{sec:future}

\begin{enumerate}[nosep]
  \item Add security assumptions made in the scheme (LRSW/PS unforgeability,
        DDH in $\GG{1}$ and $\GG{2}$, random oracles for $\Hg$, $\Hs$) and a
        proof sketch.
  \item Add the description and design of the peg-in and peg-out operations
        when joining the federation.
  \item Analyse the affine-combination attack of Section~\ref{sec:sharedh}
        enabled by the shared base $\Go{B}$ in a batched issuance, and its
        mitigations (per-note bases, range proofs at spend time, or a batching
        strategy that keeps distinct bases).
  \item Add and check the bounds provided for proof batching (soundness error
        of~\eqref{eq:batchverify}).
  \item Decide whether $\Gt{V}$ and the $\K{r'}$ terms of~\eqref{eq:rerand}
        are needed (Section~\ref{sec:V}).
  \item Fix the prototype deviations of Section~\ref{sec:impl}.
  \item Add acknowledgments.
\end{enumerate}

\begin{thebibliography}{9}
\bibitem{coconut}
A.~Sonnino, M.~Al-Bassam, S.~Bano, S.~Meiklejohn, and G.~Danezis,
``Coconut: Threshold issuance selective disclosure credentials with applications
to distributed ledgers,'' in \emph{Proceedings of the Network and Distributed
System Security Symposium (NDSS)}, Internet Society, 2019.
\url{https://arxiv.org/abs/1802.07344}.

\bibitem{nym}
Nym Technologies SA, ``The Nym network.''
\url{https://nymtech.net/nym-whitepaper.pdf}, 2021. Accessed: 2026-07-28.

\bibitem{nymcode}
Nym Technologies SA, ``nym\_offline\_compact\_ecash,'' Rust crate in the
\texttt{nymtech/nym} repository (\texttt{common/nym\_offline\_compact\_ecash}),
2024. Accessed: 2026-09-16.

\bibitem{ps16}
D.~Pointcheval and O.~Sanders, ``Short randomizable signatures,'' in
\emph{CT-RSA 2016}, LNCS 9610, Springer, 2016.

\bibitem{bulletproofs}
B.~B\"unz, J.~Bootle, D.~Boneh, A.~Poelstra, P.~Wuille, and G.~Maxwell,
``Bulletproofs: Short proofs for confidential transactions and more,'' in
\emph{IEEE Symposium on Security and Privacy}, 2018.
\end{thebibliography}

\end{document}
